\documentclass{article}
\usepackage{/globals/code}
\usepackage{/globals/anki}
\ankiprefix{methoden}
\ankitopic{Methoden}
\ankishow
\begin{document}
\section{Methoden} 
Eine Methode ist ein Abschnitt eines Programms welcher von anderswo aufgerufen werden kann.

Es kann zwischen Prozeduren (geben keinen Wert zurück) und Funktionen (geben einen Wert zurück) unterschieden werden.
\ankinote{unt}{Unterschiede}{Prozeduren (ohne Rückgabewert), Funktionen (Rückgabewert)}

\subsection{Deklaration}
Methoden haben einen Namen, einen Rückgabewert (oder \ilsourcecode{void}), und eine Anzahl an Parametern. Diese werden alle im Methodenkopf definiert. Darauf folgt der Methodenrumpf, mit den Anweisungen der Methode.
\begin{sourcecode}{java}
public static RückgabeDatentyp methodenName(ParameterDatentyp parameterName) {
    Anweisungen;
}
\ankinote{deklaration}{Deklaration}{public static RückgabeDatentyp methodenName(ParameterDatentyp parameterName)}
\end{sourcecode}
Als syntaktischer Zucker können Varargs genutzt werden; der Datentyp des letzten Parameters kann von drei Punkten gefolgt sein damit mehrere Werte des Typs eingesetzt und als Array übergeben werden können.
\begin{sourcecode}{java}
public static RückgabeDatentyp methodenName(ParameterDatentyp... parameterName) {
    Anweisungen;
}
\end{sourcecode}
\ankinote{varargs}{Varargs}{ParameterDatentyp...}

\subsection{\texttt{return;}}
Mit der \ilsourcecode{return;} Anweisung kann eine Methode verlassen werden.

Mit \ilsourcecode{return wert;} kann dabei noch ein Wert zurückgegeben werden.

\subsection{Aufruf}
Methoden können aufgerufen werden indem deren Name, gefolgt von runden Klammern, geschrieben wird.
\begin{slsourcecode}{java}
methodenName();
\end{slsourcecode}
Ist es eine Methode einer anderen Klasse, so wird der Name der Klasse vorgestellt, mit einem Punkt getrennt von dem Namen der Methode.
\begin{slsourcecode}{java}
KlassenName.methodenName();
\end{slsourcecode}
Parameter werden Kommagetrennt in die Klammern geschrieben.
\begin{slsourcecode}{java}
methodenName(wert1, wert2);
\end{slsourcecode}
\end{document}